\section*{Capítulo \ref{ch:b2:deep-time-climate} --- El clima a lo largo del tiempo geológico}

\begin{solution}{exo:b2:deep-time-climate:1}
(a) $900 + 37100 + 700 + 3 = \qty{38703}{PgC}$, unas 65 veces los \qty{591}{PgC} del aire.
(b) $\num{3.5e6}\times0.012011 = \qty{42000}{PgC}$. (c) $591/42000 = \qty{1.4}{\percent}$.
\end{solution}

\begin{solution}{exo:b2:deep-time-climate:2}
(a) Entradas $11.3 + 3 + 12.3 + 8.1 + 2.0 = \qty{36.7}{Tmol.yr^{-1}}$; salidas $26 + 16 =
\qty{42}{Tmol.yr^{-1}}$, iguales dentro de sus incertidumbres. (b) $(11.3 + 3)/36.7 = 0.39$: bastante
menos de la mitad; la mayor parte de la entrada es carbono sedimentario antiguo. (c) $\num{3.5e6}/42 = \qty{8.3e4}{yr}$.
\end{solution}

\begin{solution}{exo:b2:deep-time-climate:3}
(a) \ce{Mg2SiO4 + 4 CO2 + 2 H2O -> 2 Mg^{2+} + 4 HCO3^- + SiO2}. (b) Con \ce{2 Mg^{2+} + 4 HCO3^-
-> 2 MgCO3 + 2 CO2 + 2 H2O}, la reacción neta es \ce{Mg2SiO4 + 2 CO2 -> 2 MgCO3 + SiO2}: dos moles
de dióxido de carbono por mol de forsterita. (c) $\qty{1}{PgC} = 10^{15}/12.011 =
\qty{8.33e13}{mol}$ de carbono, que necesita \qty{4.16e13}{mol} de forsterita, $4.16\times10^{13}
\times\qty{140.7}{g} = \qty{5.86}{Pg}$, casi seis mil millones de toneladas.
\end{solution}

\begin{solution}{exo:b2:deep-time-climate:4}
Glaciación hirnantiense, de 445 a \qty{443}{Ma}, estado glaciar; máximo del hielo del Paleozoico superior, de 307 a
\qty{299}{Ma}, estado glaciar; Turoniense, de 94 a \qty{89}{Ma}, invernadero (el más cálido del Fanerozoico);
PETM, de 55.9 a \qty{55.7}{Ma}, invernadero; EECO, de 53 a \qty{49}{Ma}, invernadero; inicio del casquete de hielo
antártico, hacia \qty{34}{Ma}, comienzo del estado glaciar actual; casquetes de hielo del norte,
\qty{2.7}{Ma}, estado glaciar.
\end{solution}

\begin{solution}{exo:b2:deep-time-climate:5}
(a) $63000/(8.314\times288^2) = \qty{0.0914}{K^{-1}}$, que redondeado da 0.09. (b) $74000/(8.314
\times288^2) = \qty{0.107}{K^{-1}}$. (c) $\ln1.2/0.09 = \qty{2.03}{K}$ y $\ln 1.2/0.107 =
\qty{1.70}{K}$. (d) En el estado estacionario la meteorización debe igualar a la desgasificación, lo que fija $e^{k
\Delta T}$ y por tanto $\Delta T$; la sensibilidad solo determina cuánto dióxido de carbono produce ese
calentamiento.
\end{solution}

\begin{solution}{exo:b2:deep-time-climate:6}
Aquí $k_T = 0.0914$. (a) Exacto $\exp[7577.5\times(1/288 - 1/293)] = 1.567$; aproximación
$e^{0.0914\times5} = 1.579$: un \qty{0.8}{\percent} por exceso. (b) Exacto $e^{2.482} = 11.97$;
aproximación $e^{2.741} = 15.50$: un \qty{30}{\percent} por exceso. (c) $\ln(\text{razón}) = (E_{\text
a}/R)[x/T_0^2 - x/(T_0(T_0 + x))] = (E_{\text a}/R)\,x^2/(T_0^2(T_0 + x))$. Igualándolo a $\ln1.1$:
$x^2 = 1.043\,(288 + x)$, así que $x = \qty{17.9}{K}$. Más allá de unos \qty{18}{K} de calentamiento, la
forma exponencial sobrestima la meteorización en más de una décima parte.
\end{solution}

\begin{solution}{exo:b2:deep-time-climate:7}
(a) $1444/276 = 5.23$, y $3.93\log_2 5.23 = \qty{9.4}{W.m^{-2}}$. (b) Casi exactamente los
\qty{9}{W.m^{-2}} del Sol: los dos forzamientos se compensan. (c) $(1444 - 276)/419.5 =
\qty{2.8}{ppm}$ por millón de años. Los \qty{3.4}{ppm} proceden de una recta ajustada a todos los
datos, que pondera cada estimación en lugar de los dos extremos de la curva suavizada.
\end{solution}

\begin{solution}{exo:b2:deep-time-climate:8}
(a) De $3000/20000 = \qty{0.15}{PgC.yr^{-1}}$ a $7000/3000 = \qty{2.3}{PgC.yr^{-1}}$, un intervalo
que contiene el de 0.3 a 1.5 del IPCC. (b) $35\times12.011/44.01 = \qty{9.6}{PgC.yr^{-1}}$, de 6 a
32 veces el ritmo de 0.3 a \qty{1.5}{PgC.yr^{-1}} del PETM. (c) $3.93\log_2(2000/900) = \qty{4.5}{W.m^{-2}}$,
y $3\log_2 2.22 = \qty{3.5}{\celsius}$: alrededor del setenta por ciento de los \qty{5}{\celsius}. Otras
retroalimentaciones u otra fuente de calentamiento, como el metano, deben aportar el resto.
\end{solution}

\begin{solution}{exo:b2:deep-time-climate:9}
(a) Ahora $1.2\,W_0e^{k\Delta T} = V = W_0$, así que $\Delta T = -\ln 1.2/0.09 = \qty{-2.03}{K}$.
(b) $278\times2^{-2.03/3} = \qty{174}{ppm}$. (c) No: en cualquier estado estacionario la meteorización iguala a la
desgasificación, que no ha cambiado. La mayor meteorizabilidad queda compensada por un clima más frío.
\end{solution}

\begin{solution}{exo:b2:deep-time-climate:10}
(a) $k = 0.119$: $\tau = 591\times0.693/(0.08\times0.119\times3\times0.138) = \qty{104}{kyr}$.
(b) $S = 8$: $\tau = \qty{51.5}{kyr}$. (c) Ambos hacen más rápido el termostato, ya que $\tau \propto
1/(kS)$. $E_{\text a}$ es incierta en un factor dos aproximadamente, y $S$, a estas escalas de tiempo, en un
factor parecido (3 frente a 5--\qty{8}{K}): ambas pesan más o menos lo mismo.
\end{solution}

\begin{solution}{exo:b2:deep-time-climate:11}
(a) $0.08\times5000 = \qty{400}{PgC}$ permanecen en el aire: $C/C_0 = 991/591 = 1.677$, es decir,
\qty{466}{ppm}; $\Delta T = 3\log_2 1.677 = \qty{2.24}{K}$. (b) $W - W_0 = 0.138\,(e^{0.09\times2.24}
- 1) = \qty{0.0308}{PgC.yr^{-1}}$, un \qty{22}{\percent} más de meteorización; $5000/0.0308 =
\qty{162}{kyr}$. (c) $\tau = \qty{137}{kyr}$ y $\tau\ln10 = \qty{315}{kyr}$. (d) El calentamiento crece
solo como el logaritmo del dióxido de carbono, así que un exceso grande aumenta la meteorización menos de lo que
supone la ley lineal; la retirada empieza más despacio.
\end{solution}

\begin{solution}{exo:b2:deep-time-climate:12}
(a) Cerca de $M^*$, $\mathrm d\Delta T/\mathrm dM = Sf/(A_0\ln 2\,(1 + fM^*/A_0)) = Sf/(A_0\ln 2\,
C^*/C_0)$ y $\mathrm dW/\mathrm dM = kW\,\mathrm d\Delta T/\mathrm dM$, con $W = V$ allí; la
inversa es $\tau'$. (b) $C^*/C_0 = 2^{2.03/3} = 1.597$: $\tau' = 137.4\times1.597/1.2 =
\qty{183}{kyr}$. (c) Más lento: con mucho dióxido de carbono, cada petagramo de carbono añadido calienta menos,
porque el forzamiento es logarítmico, así que la meteorización responde menos a él.
\end{solution}

\begin{solution}{pb:b2:deep-time-climate:1}
\textbf{1.} $1.3 + 1.3 + 4.5 + 4.2 = \qty{11.3}{Tmol.yr^{-1}}$, $\qty{0.136}{PgC.yr^{-1}}$.
\textbf{2.} $11.5\times0.012011 = \qty{0.138}{PgC.yr^{-1}}$.
\textbf{3.} Coinciden con un margen de un \qty{2}{\percent} aproximadamente: la meteorización de silicatos equilibra la desgasificación
volcánica, como exige el \omterm{def:b2:deep-time-climate:carbon-cycle}{ciclo del carbono a largo plazo}.
\textbf{4.} $\num{5.1e18}/0.02897 = \qty{1.76e20}{mol}$ de aire; $\num{591e15}/12.011 =
\qty{4.92e16}{mol}$ de carbono; $\num{4.92e16}/\num{1.76e20} = \num{2.80e-4}$, \qty{280}{ppm}.
\textbf{5.} $591/0.136 = \qty{4350}{yr}$.
\textbf{6.} $1 + 0.08 = 1.08$: solo un \qty{8}{\percent} más.
\textbf{7.} $63000/(8.314\times288^2) = \qty{0.0914}{K^{-1}}$, que redondeado da 0.09.
\textbf{8.} $e^{0.09} = 1.09$: un \qty{9}{\percent} más por kelvin.
\textbf{9.} $\Delta T = S\log_2(1 + fM/A_0)$; su derivada en 0 es $Sf/(A_0\ln2) =
3\times0.08/(591\times0.693) = \qty{5.86e-4}{K.PgC^{-1}}$.
\textbf{10.} $W_0k\times\num{5.86e-4} = 0.138\times0.09\times\num{5.86e-4} = \qty{7.28e-6}{yr^{-1}}$.
\textbf{11.} Su inversa es $A_0\ln2/(fkSW_0)$, la $\tau$ de la proposición, ya que la
ecuación linealizada es $\mathrm dM/\mathrm dt = -(\mathrm dW/\mathrm dM)\,M$.
\textbf{12.} Los cuatro: $\tau$ es inversamente proporcional a cada uno.
\textbf{13.} $1/\num{7.28e-6} = \qty{1.37e5}{yr}$.
\textbf{14.} $137\times0.09/0.119 = \qty{104}{kyr}$.
\textbf{15.} $k_T = 0.107$: $137\times0.09/0.107 = \qty{115}{kyr}$.
\textbf{16.} Más cortos en un factor de 1.5 a 7: el modelo de una caja y su $f$ constante hacen el
termostato demasiado rápido.
\textbf{17.} $\num{50e6}/\num{1.37e5} = 364$: el termostato tuvo tiempo de equilibrar el ciclo en cada
etapa, así que el descenso siguió cambios lentos de la desgasificación o de la meteorizabilidad.
\textbf{18.} $16\times0.137 = \qty{2.2}{ppm}$.
\textbf{19.} $0.08\times5000 = \qty{400}{PgC}$; $278\times991/591 = \qty{466}{ppm}$.
\textbf{20.} $3\log_2 1.677 = \qty{2.24}{K}$.
\textbf{21.} $0.138\,(e^{0.202} - 1) = \qty{0.0308}{PgC.yr^{-1}}$, un \qty{22.3}{\percent}.
\textbf{22.} $137\times\ln10 = \qty{315}{kyr}$.
\textbf{23.} La meteorización pudo ser más sensible (escorrentía, una energía de activación mayor, las
retroalimentaciones lentas de una mayor sensibilidad del sistema Tierra); y el carbono orgánico en tierra y en el mar pudo
recuperarse y retirar carbono del aire más deprisa que la meteorización de silicatos.
\textbf{24.} En este modelo, el \omterm{def:b2:deep-time-climate:thermostat}{termostato de la meteorización de los silicatos} retira un exceso de carbono con un
tiempo característico de unos \textbf{\qty[text-series-to-math=true]{137}{} mil años}.
\end{solution}
